\section{Discussion and Limitations}
\label{sec:discussion}

\emph{Scale.} The main evaluation covers \E{full:prompts} prompts from three task families, with one observation per prompt, method and delay. The prompt-cluster intervals account for prompt-to-prompt variation, and the held-out sets and the subset of previously unused prompts agree with the main result. Repeated runs would additionally capture run-to-run variation of the cloud and the network.

\emph{Generality.} We study one device, one model pair and one server GPU. \cape{} measures its cost inputs ($V$, $d$) and acceptance rates online, so we expect it to transfer to other model pairs and devices with only the power ratio $r$ recalibrated, but we have not yet tested this.

\emph{Network conditions.} All results include the jitter of a real campus Wi-Fi path with both endpoints on wireless links, whose per-request overhead exceeded 100\,ms in \E{net:over100} of rounds; the median-of-three round-trip estimate absorbs such spikes. The added delays, however, are constant, and we have not evaluated cellular access networks or wide-area traces.

\emph{Energy scope and trade-off.} We report gross board energy of the edge device, which includes idle power. The energy of the RTX 3080 server and of the network is not measured, so our results concern the device's energy budget, not the system's. A data-center GPU would also verify faster than a desktop RTX 3080, which shortens $V$ and favors shorter lookahead; \cape{} adapts to this because it measures $V$ online. The speed setting \capes{} uses \E{full:capes:fixed4:pooled:abssv} more J/token than Fixed-4 overall and \E{full:capes:fixed4:high:abssv} more at 100\,ms. Applications that must not spend more energy per token should use $w\ge0.25$, which gave speedups without a measurable energy penalty on the held-out sets, or $w=0.5$, which saves energy at unchanged speed.

\emph{Calibration of $r$.} The controller uses $r=2.3$ from development runs. Fitting Eq.~\eqref{eq:energy} to our data sets gives $r$ between \E{energy:rmin} and \E{energy:rmax}. Wall time and draft-busy time are correlated across requests, so the split between the two coefficients is less certain than their total, and the calibrated value lies within the fitted range. Because $r$ only scales the energy term, a mis-calibrated $r$ acts like a different energy weight. Online calibration from the power sensor is a natural extension.

\emph{Exactness.} Strict greedy verification makes every emitted token the target's top choice, but with a quantized target near-ties resolve differently for different verification shapes. Outputs therefore differ between methods for some prompts, as they do between fixed windows of different sizes and between speculative and autoregressive decoding (Section~\ref{sec:integrity}). GSM8K accuracy was unaffected. HumanEval pass@1 and MT-Bench quality scores were not evaluated.

\emph{Cloud-side opportunity.} Verification time is dominated by recomputing the context, because requests for prompt log-probabilities bypass the prefix cache in our deployment. Incremental verification with cached prefixes would shrink $V$ substantially. This is complementary to \cape{}, which would respond by shortening its lookahead, since it measures $V$ online.
